\documentclass[11pt]{article}
\usepackage[margin=1in]{geometry}
\usepackage[T1]{fontenc}
\usepackage{lmodern,microtype,amsmath,amssymb,amsthm,graphicx,booktabs,tabularx}
\usepackage[colorlinks=true,linkcolor=blue,urlcolor=blue]{hyperref}
\setlength{\parindent}{0pt}
\setlength{\parskip}{4pt}
\newtheorem{proposition}{Proposition}
\newcommand{\dd}{\mathrm d}
\begin{document}
\begin{center}
{\Large Housing allocation, fertility, and population}\par
\smallskip
{\small Proposed results following the existing household setup}
\end{center}

\section*{Housing allocation}

Borrowing limits can leave housing with old owners even when young households
would give up more consumption to use it. In an equilibrium of the model,
write $x_i=c_i-\chi n_i$ and $s_i=h_i-\kappa n_i$ for
the goods and space available to young adults. The young household's and old
owner's marginal values of housing, in current consumption goods, are:
\begin{equation}
 MV_i^Y=\frac{\alpha x_i}{s_i},\qquad
 MV_j^O=\frac{\gamma c_j^2}{h_j^2}.
 \label{eq:core_values}
\end{equation}
These values measure what each household would pay for a little more space.

The planner reallocates housing and consumption, retaining the physical
limits on each tenure and the total housing stock. It may relax private
financing restrictions. Each household's fertility, tenure, estate, and
future real allocation remain fixed, and existing financial claims are honored.

\begin{proposition}[Housing misallocation]
Suppose positive masses of young owners can occupy more housing and old
owners can release some. If $MV_i^Y>MV_j^O$ on matched groups, a sufficiently
small reallocation toward the young can make them better off and leave
every other household equally well off, in the absence of real reallocation costs.
\end{proposition}

\textit{Proof.} Give the young household $\epsilon$ units of the old owner's
housing. Keeping the old owner's utility fixed requires extra consumption:
\begin{equation}
 D_j(\epsilon)=c_j^2\left[
 \left(\frac{h_j^2}{h_j^2-\epsilon}\right)^\gamma-1\right].
 \label{eq:core_compensation}
\end{equation}
Take this amount from the young household. Goods and housing balance, and
its utility gain has derivative $(MV_i^Y-MV_j^O)/x_i>0$ at zero.
A small change preserves feasibility.
Financial settlement replaces the seller's lost estate title with an equal
bond claim; the buyer's later sale of the extra title offsets its financing.
Future allocations, existing creditors, and aggregate external assets are
unchanged.\hfill$\square$

\begin{center}
\includegraphics[width=.70\textwidth]{output/model/simplified_olg_amendments/simple_allocation_figure.pdf}
\end{center}
{\small\textbf{Figure 1.} Housing moves toward the young while the old
owner is compensated. The shaded gap indicates the direction of improvement;
its area is not a welfare measure. Equality of values gives the best allocation
along this particular compensated comparison.}

\newpage
\section*{Housing access and fertility}

The borrowing limit can produce the inequality in Proposition 1. With
positive saving, slack physical owner limits, and interior old-age choices,
a strictly restrictive down payment gives $MV_i^Y>q r_t$. An old owner with
slack retention and estate bounds has $MV_j^O=q r_t$. Young buyers need cash
at purchase; old owners already hold their homes.

Rental segmentation limits another way to obtain space. A household may be
unable to rent a larger home and unable to fund its purchase. The ownership
taste in the model also matters: it can make constrained ownership preferable
even when larger rentals are available. Thus both access restrictions can
coexist with misallocation, although a binding rental cap is not necessary
for the owner-to-owner result. A renter already at a physical cap cannot
receive more housing in that tenure under the same planner restrictions.

Real reallocation costs have a simple implication. If moving $\epsilon$
housing uses $C(\epsilon)$ goods, with $C(0)=0$ and $C'(0)=K_{ijt}\geq0$, the condition
becomes $MV_i^Y>MV_j^O+K_{ijt}$. A positive fixed cost instead requires a finite
reallocation whose gain covers that cost.

\textbf{A fertility prediction.} More space makes children easier to
accommodate; paying for it reduces the goods available to the family.
Consider the young owner's conditional problem with positive saving and
interior old-age choices. Suppress household and date indices, and let
$k=\beta(1+\gamma+\omega_B)$. The fertility condition and the response of
the lifetime budget to more housing are:
\begin{equation}
 \frac{\vartheta}{n}=\frac{\chi}{x}+\frac{\alpha\kappa}{s},
 \qquad (1+k)\dd x=-\chi\dd n-p\dd h.
 \label{eq:core_fertility_system}
\end{equation}
Here $p$ is the net lifetime payment for a marginal unit of current housing,
including accompanying transfers. With prices and other resources fixed,
$p=q r_t$. Differentiating the first equation using the second gives:
\begin{equation}
 \frac{\dd n}{\dd h}>0
 \quad\Longleftrightarrow\quad
 \frac{\alpha\kappa}{s^2}>
 \frac{\chi p}{(1+k)x^2}.
 \label{eq:core_fertility_sign}
\end{equation}
Fertility rises when the space benefit exceeds the resource cost. This is a
household response while tenure and the binding constraints remain unchanged.
Saving and old-age choices are reoptimized. Proposition 1 instead holds
future real allocations fixed.

A sufficient condition can be written without prices or a borrowing
multiplier. At stationary prices, with zero property tax and a strictly
binding down-payment limit, it is:
\begin{equation}
 \frac{(1-q)b}{(1-\phi)(y+b)}
 \geq\frac{\alpha}{1+\beta(1+\gamma+\omega_B)+\alpha}.
 \label{eq:core_primitive}
\end{equation}
The left side is housing expenditure at the purchase limit relative to
lifetime resources. Under the stated household conditions, this restriction
ensures a fertility increase for any $p<MV^Y$. The full equilibrium response
also depends on prices, tenure choices, and subsequent cohort sizes.

\newpage
\section*{Fertility and population over time}

A temporary fertility increase can leave a permanently larger population.
Let $Y_t$ and $O_t$ count young and old households, $\bar n_t$ denote mean
young fertility, and $\nu$ convert children into entering households. Two
paths starting from the same young cohort satisfy:
\begin{equation}
 \frac{Y_t^1}{Y_t^0}=
 \prod_{s=0}^{t-1}\frac{\bar n_s^1}{\bar n_s^0},
 \qquad O_t^j=Y_{t-1}^j.
 \label{eq:core_population}
\end{equation}
Population reflects cumulative fertility differences. Both positive steady
states have replacement fertility $\bar n^*=1/\nu$. Housing clearing gives:
\begin{equation}
 N_{\rm hh}^*=2Y^*=
 \frac{2\bar H}{\bar h^{Y*}+\bar h^{O*}}.
 \label{eq:core_stationary_population}
\end{equation}
With a fixed stock, a larger population requires less housing per household
on average. This stationary ranking also applies to persons under unchanged
household and child-counting conventions.

The model admits local equilibrium transitions in which a small permanent
increase in the financed share raises initial fertility and final population.
This result holds in a family of economies with almost half of households
renting, positive child goods costs, and a positive rebated property tax.
Tenure choice remains endogenous. Small changes in the primitives and small
entrant income and wealth heterogeneity preserve the result, provided the
same household constraints continue to bind.\footnote{The accompanying
mixed-tenure appendix gives the primitive family and proof. Young owners
remain at their purchase limits, renters at their physical caps, and saving
is positive in both tenures. Old owners' retention and estate bounds are
slack. The reform must be sufficiently small; no numerical reform-size bound
is established. The proof combines exact algebra and rational interval
bounds, independently checked against the original household equations.}

Initial fertility and final population are separate comparisons. A larger
final population depends on cumulative fertility, and does not require
fertility to exceed replacement at every date.

\begin{center}
\includegraphics[width=\textwidth]{output/model/simplified_olg_amendments/mixed_transition_analytical_figure.pdf}
\end{center}
{\small\textbf{Figure 2.} Analytical first-order equilibrium response in
the mixed-tenure economy, per unit increase in $\phi$. Initial fertility
rises and population ends higher. The damped oscillations also hold for
sufficiently small finite reforms near this economy: fertility approaches
replacement from both sides.}

In the illustrated economy, part of the increased borrowing capacity enters
house prices. Conditional young owners buy larger homes, old owners use less
housing, and ownership becomes less common. Average housing per household
falls, allowing the fixed stock to accommodate more households. The welfare
gain in Proposition 1 comes from a separately compensated reallocation.

\newpage
\section*{For discussion: how much theory to retain}

\textbf{The allocation result can carry the main claim.} It identifies a
feasible welfare gain, explains why housing can remain with the old, and
requires one short proof. The planner's borrowing permission is the one
already accepted. A complete first-best allocation and a market-based
implementation remain separate questions.

\textbf{The transition result covers a mixed economy, but remains local.}
It no longer requires almost universal ownership or a child goods cost close
to zero. The following economic restrictions still matter:

\renewcommand{\arraystretch}{1.2}
\begin{tabularx}{\textwidth}{@{}p{.29\textwidth}>{\raggedright\arraybackslash}X@{}}
\toprule
Assumption & Economic consequence\\
\midrule
Small reforms and unchanged binding constraints
& The proof permits tenure switching but does not follow large reforms across
changes in the conditional housing or saving constraints.\\
Old owners can sell; estate bounds are slack
& The interior stationary formula requires slack retention and
\mbox{$[1-q+q\tau^p]\omega_B>q\gamma$}. A binding bound changes old housing demand.\\
Fixed entrant endowments and world bond price
& Larger cohorts receive the same distribution of income and wealth. Wages,
interest rates, and inherited wealth do not adjust to the larger population.\\
\bottomrule
\end{tabularx}

A small reform is a reasonable scope for an illustrative result. The fixed
income and wealth of new entrants are a stronger limitation. The construction
establishes a class of economies with the desired demographic response; it
does not establish how common that response is empirically. The earlier
all-owner example and its simpler primitive conditions remain available in
the longer assessment.

\textbf{The demographic gain is not a welfare ranking.}
Throughout the constructed mixed-tenure family, the stationary comparison
has the following signs, holding both ownership-taste parameters fixed:
\begin{equation}
 \frac{\dd Y^*}{\dd\phi}>0,\qquad
 \frac{\dd W^{O*}}{\dd\phi}<0,\qquad
 \frac{\dd W^{R*}}{\dd\phi}<0.
 \label{eq:core_stationary_welfare}
\end{equation}
Both conditional lifetime values fall, so the same entrant is worse off
regardless of whether its fixed taste draw leads it to switch tenure. This
compares that entrant in two steady states. It does not assign a social value
to the change in population or evaluate every transitional cohort. Initial
old owners may benefit from the appreciation of their inherited titles.

\textbf{My recommendation.} Use the allocation result and the conditional
fertility condition as the core. Keep both figures available, presenting the
transition as a separate illustrative extension and putting its proof in the
appendix. Its value is the connection between temporary fertility and lasting
population, with the price response made explicit. Whether that extension
earns main-text space is the next useful discussion. Constrained efficiency
and public lending can remain later branches. No model convention or planner
power is changed by this proposed presentation.

\end{document}
